\chapter{Nonlinear projection-based model order reduction using machine learning regression to model closure error in the latent space (Article from International Research Stay)}
\label{chap:article4}
This chapter presents the article titled \emph{``Nonlinear projection-based model order reduction using machine learning regression to model closure error in the latent space''}~\cite{de2025nonlinear}, published in \textit{Computer Methods in Applied Mechanics and Engineering} (CMAME). This work was carried out as part of a long-term Fulbright predoctoral research stay as a Visiting Researcher at Stanford University within the research group of Professor Charbel Farhat.

Building on the latent-space closure modeling framework introduced in earlier chapters, this contribution investigates alternative, more interpretable regression strategies that address key limitations of PROM-ANN, such as its reliance on black-box neural networks and large training datasets. Specifically, it proposes Gaussian Process Regression (PROM-GPR) and Radial Basis Function interpolation (PROM-RBF) as kernel-based surrogates for modeling closure errors in the latent space, preserving the flexibility of the nonlinear manifold while enhancing interpretability, data efficiency, and potential certification.

The paper demonstrates:
\begin{itemize}
    \item A unified formulation of PROM-GPR and PROM-RBF within the general nonlinear manifold ROM framework, including analytical expressions for the kernel-based closure maps and their derivatives.
    \item The integration of these regression surrogates into the residual-minimizing Petrov--Galerkin projection and their compatibility with project-then-approximate hyper-reduction techniques.
    \item Validation on industrial-like cases, including a detached-eddy simulation of the Ahmed body wake flow and the two-dimensional inviscid Burgers' equation, illustrating how the approach delivers accuracy comparable to PROM-ANN with significantly reduced training data requirements.
\end{itemize}

This study showcases how kernel-based regression strategies, coupled with robust projection and hyper-reduction frameworks, can effectively mitigate the effects of the Kolmogorov $n$-width barrier in convection-dominated and strongly nonlinear systems. It represents a key contribution of this thesis toward scalable, interpretable, and industry-relevant nonlinear PROMs.

\includepdf[
  pages=-,
  pagecommand={},
  scale=0.92,
  delta=0 0,
  offset=32pt 0pt,
  turn=false
]{Papers/Nonlinear_projection_based_model_order_reduction_with_machine_learning_regression_for_closure_error_modeling_in_the_latent_space.pdf}

